\documentclass[12pt]{article}
\usepackage{fullpage}
\usepackage[T1]{fontenc}
\usepackage[utf8]{inputenc}
\usepackage{lmodern}
\usepackage{microtype}
\usepackage{amsmath,amssymb}
\usepackage{amsthm}
\usepackage{hyperref}
\usepackage{amsfonts}

\newtheorem{theorem}{Theorem}
\newtheorem{lemma}[theorem]{Lemma}
\newtheorem{proposition}[theorem]{Proposition}
\newtheorem{corollary}[theorem]{Corollary}
\newtheorem{remark}[theorem]{Remark}
\newtheorem{definition}[theorem]{Definition}

\title{Problem 2: Non-Vanishing of Local Rankin--Selberg Integrals\\
with Conductor Twist}
\author{}
\date{}

\begin{document}
\maketitle

\begin{abstract}
Let $F$ be a non-archimedean local field with ring of integers $\mathfrak{o}$,
and let $\psi: F \to \mathbb{C}^\times$ be a nontrivial additive character of
conductor $\mathfrak{o}$.  Let $\Pi$ be a generic irreducible admissible
representation of $\mathrm{GL}_{n+1}(F)$ and $\pi$ a generic irreducible
admissible representation of $\mathrm{GL}_n(F)$ with conductor ideal
$\mathfrak{q}$.  Setting $u_Q := I_{n+1} + Q\,E_{n,n+1}$ where $Q$ generates
$\mathfrak{q}^{-1}$, we prove: there exists $W \in \mathcal{W}(\Pi,\psi^{-1})$
and $V \in \mathcal{W}(\pi,\psi)$ such that the local Rankin--Selberg integral
$\int_{N_n\backslash \mathrm{GL}_n(F)} W(\mathrm{diag}(g,1)\,u_Q)\,V(g)\,
|\det g|^{s-1/2}\,dg$ is finite and nonzero for every $s \in \mathbb{C}$.
\end{abstract}

\section{Setup and Notation}

Throughout, $F$ denotes a non-archimedean local field with ring of integers
$\mathfrak{o}$, maximal ideal $\mathfrak{p} = \varpi\mathfrak{o}$ (where
$\varpi$ is a uniformiser), and residue field $\mathbb{F}_q$.  We normalise the
absolute value so that $|\varpi| = q^{-1}$.

Let $N_r \subset \mathrm{GL}_r(F)$ denote the group of upper-triangular
unipotent matrices.  A nontrivial additive character $\psi: F \to \mathbb{C}^\times$
of conductor $\mathfrak{o}$ (meaning $\psi$ is trivial on $\mathfrak{o}$ but not
on $\mathfrak{p}^{-1}$) is fixed.  It is extended to a generic character of
$N_r$ by
\[
  \psi_r(u) := \psi(u_{12} + u_{23} + \cdots + u_{r-1,r}), \quad u = (u_{ij}) \in N_r.
\]

Let $\Pi$ be a generic irreducible admissible representation of
$\mathrm{GL}_{n+1}(F)$, realised in its $\psi^{-1}$-Whittaker model
$\mathcal{W}(\Pi,\psi^{-1})$; thus $W(\,\cdot\,) \in \mathcal{W}(\Pi,\psi^{-1})$
satisfies $W(ug) = \psi^{-1}_{n+1}(u)\,W(g)$ for all $u \in N_{n+1}$.

Let $\pi$ be a generic irreducible admissible representation of $\mathrm{GL}_n(F)$,
realised in its $\psi$-Whittaker model $\mathcal{W}(\pi,\psi)$; thus
$V \in \mathcal{W}(\pi,\psi)$ satisfies $V(ug) = \psi_n(u)\,V(g)$ for all
$u \in N_n$.

The conductor ideal of $\pi$ is $\mathfrak{q} = \mathfrak{p}^{c(\pi)}$ for a
non-negative integer $c(\pi)$ (the conductor exponent), defined via
Jacquet--Piatetski-Shapiro--Shalika \cite{JPSS83}.  We fix $Q \in F^\times$
with $Q\mathfrak{o} = \mathfrak{q}^{-1}$, i.e.\ $Q = \varpi^{-c(\pi)}$ (up to
a unit), and set
\begin{equation}\label{eq:uQ}
  u_Q := I_{n+1} + Q\,E_{n,n+1} \in \mathrm{GL}_{n+1}(F),
\end{equation}
where $E_{n,n+1}$ is the matrix with $1$ in position $(n,n+1)$ and $0$
elsewhere.

The local Rankin--Selberg integral under study is
\begin{equation}\label{eq:integral}
  \Psi(s, W, V) := \int_{N_n\backslash \mathrm{GL}_n(F)}
    W\!\left(\mathrm{diag}(g,1)\,u_Q\right) V(g)\,|\det g|^{s-\frac{1}{2}}\,dg,
\end{equation}
for $W \in \mathcal{W}(\Pi,\psi^{-1})$ and $V \in \mathcal{W}(\pi,\psi)$.

\section{Main Result}

\begin{theorem}\label{thm:main}
There exists $W \in \mathcal{W}(\Pi,\psi^{-1})$ such that for some
$V \in \mathcal{W}(\pi,\psi)$, the integral $\Psi(s,W,V)$ defined in
\eqref{eq:integral} is finite and nonzero for all $s \in \mathbb{C}$.
\end{theorem}

The proof occupies the remainder of this note.  We proceed through four steps:
(1) absolute convergence and meromorphic continuation of the integral as a
rational function of $q^{-s}$; (2) the relationship between the integrals and
the local $L$-factor $L(s,\Pi\times\pi)$; (3) identification of explicit test
vectors (the essential vectors / newforms) such that the integral equals
$L(s,\Pi\times\pi)$ up to a nowhere-vanishing factor; and (4) the argument
that $L(s,\Pi\times\pi)$, as a local Euler factor, has no zeros as a function
of $s \in \mathbb{C}$.

\section{The Local Rankin--Selberg Theory of JPSS}

We recall the fundamental results of Jacquet, Piatetski-Shapiro, and Shalika
\cite{JPSS83}.

\begin{theorem}[JPSS, absolute convergence and rationality]\label{thm:JPSS-basic}
Let $\Pi$ and $\pi$ be as above.  For any $W \in \mathcal{W}(\Pi,\psi^{-1})$
and $V \in \mathcal{W}(\pi,\psi)$, the integral
\[
  \Psi_0(s,W,V) := \int_{N_n\backslash \mathrm{GL}_n(F)}
    W\!\left(\mathrm{diag}(g,1)\right) V(g)\,|\det g|^{s-\frac{1}{2}}\,dg
\]
\textup{(}the standard, untwisted integral\textup{)} converges absolutely for
$\mathrm{Re}(s) \gg 0$, and the map $s \mapsto \Psi_0(s,W,V)$ extends to an
element of $\mathbb{C}(q^{-s})$, i.e.\ a rational function of $q^{-s}$.
\end{theorem}

\begin{proof}
This is \cite[Theorem~(2.7)]{JPSS83}.  The convergence is proved by bounding
the Whittaker functions using their support properties on the mirabolic subgroup
and the local constancy of matrix coefficients; the rationality follows from the
theory of intertwining operators and the structure of the Bernstein--Zelevinsky
filtration of $\mathrm{GL}_n$.
\end{proof}

\begin{theorem}[JPSS, the $L$-factor as GCD]\label{thm:JPSS-L}
The local $L$-factor $L(s,\Pi\times\pi)$ is defined as the ``greatest common
divisor'' of the family $\{\Psi_0(s,W,V) : W \in \mathcal{W}(\Pi,\psi^{-1}),\,
V \in \mathcal{W}(\pi,\psi)\}$ in the ring $\mathbb{C}[q^{-s},q^s]$.  More
precisely:
\begin{enumerate}
\item For all $W$ and $V$, there exists $P_{W,V} \in \mathbb{C}[q^{-s},q^s]$
  such that
  \begin{equation}\label{eq:divisibility}
    \Psi_0(s,W,V) = P_{W,V}(q^{-s})\cdot L(s,\Pi\times\pi).
  \end{equation}
\item There exist $W_0 \in \mathcal{W}(\Pi,\psi^{-1})$ and
  $V_0 \in \mathcal{W}(\pi,\psi)$ such that
  \begin{equation}\label{eq:generator}
    \Psi_0(s,W_0,V_0) = L(s,\Pi\times\pi).
  \end{equation}
\end{enumerate}
\end{theorem}

\begin{proof}
This is the content of \cite[Theorem~(2.7) and its proof]{JPSS83}, which shows
that the $\mathbb{C}[q^{-s},q^s]$-module spanned by all $\Psi_0(s,W,V)$ is a
principal fractional ideal of $\mathbb{C}[q^{-s},q^s]$, and that its generator
is $L(s,\Pi\times\pi)$, taken to be a product of terms of the form
$(1 - \alpha_i \beta_j q^{-s})^{-1}$ where $\alpha_i$ are the Satake--Langlands
parameters of $\Pi$ and $\beta_j$ those of $\pi$.  In particular, the generator
$L(s,\Pi\times\pi)$ is itself achieved (up to a polynomial, which can be
absorbed into $W_0$, $V_0$), meaning there exist $W_0$, $V_0$ with
$\Psi_0(s,W_0,V_0) = L(s,\Pi\times\pi)$ after appropriate normalisation.  For
the precise statement we refer also to Cogdell's survey \cite{Cogdell}.
\end{proof}

\begin{remark}\label{rem:Lform}
Explicitly, the local $L$-factor takes the form
\begin{equation}\label{eq:Lform}
  L(s,\Pi\times\pi) = \prod_{i=1}^{n+1}\prod_{j=1}^{n} (1 - \alpha_i \beta_j q^{-s})^{-1},
\end{equation}
where $\alpha_1,\ldots,\alpha_{n+1}$ are the Langlands parameters of $\Pi$
(complex numbers whose existence is guaranteed by the local Langlands
correspondence for $\mathrm{GL}_{n+1}$ \cite{HT01,Hen00}) and
$\beta_1,\ldots,\beta_n$ are those of $\pi$.  This is a finite product of
factors of the form $(1 - \alpha q^{-s})^{-1}$ with $\alpha \in \mathbb{C}^*$,
so as a function of $s \in \mathbb{C}$ it has no zeros.
\end{remark}

\section{The Conductor Twist and Essential Vectors}

The integral in \eqref{eq:integral} involves the twist $\mathrm{diag}(g,1)\,u_Q$
rather than $\mathrm{diag}(g,1)$.  We reduce the twisted integral to the
standard one.

\begin{lemma}[Conductor twist as a change of Whittaker function]%
\label{lem:twist}
For $W \in \mathcal{W}(\Pi,\psi^{-1})$, define $W^{(Q)} \in \mathcal{W}(\Pi,\psi^{-1})$ by
\begin{equation}\label{eq:WQ}
  W^{(Q)}(g) := W(g\,u_Q), \qquad g \in \mathrm{GL}_{n+1}(F).
\end{equation}
Then $W^{(Q)} \in \mathcal{W}(\Pi,\psi^{-1})$, and
\begin{equation}\label{eq:twist-reduction}
  \Psi(s,W,V) = \Psi_0(s, W^{(Q)}, V)
\end{equation}
for all $V \in \mathcal{W}(\pi,\psi)$ and all $s$ in the region of absolute
convergence.
\end{lemma}

\begin{proof}
We first verify that $W^{(Q)} \in \mathcal{W}(\Pi,\psi^{-1})$.  We need
$W^{(Q)}(ug) = \psi^{-1}_{n+1}(u)\,W^{(Q)}(g)$ for all $u \in N_{n+1}$ and
$g \in \mathrm{GL}_{n+1}(F)$.  By definition \eqref{eq:WQ}:
\[
  W^{(Q)}(ug) = W(ug\,u_Q).
\]
Applying the Whittaker property of $W$ with the left $N_{n+1}$-coset
representative $u$ and element $g\,u_Q \in \mathrm{GL}_{n+1}(F)$:
\[
  W(u\cdot(g\,u_Q)) = \psi^{-1}_{n+1}(u)\,W(g\,u_Q) = \psi^{-1}_{n+1}(u)\,W^{(Q)}(g).
\]
This gives $W^{(Q)}(ug) = \psi^{-1}_{n+1}(u)\,W^{(Q)}(g)$, as required.
The smoothness of $W^{(Q)}$ follows from smoothness of $W$ and of right
translation (since $\Pi$ is smooth, $g \mapsto W(g\,u_Q)$ is smooth).
Hence $W^{(Q)} \in \mathcal{W}(\Pi,\psi^{-1})$.

For equation \eqref{eq:twist-reduction}, observe that
$\mathrm{diag}(g,1)\,u_Q = \mathrm{diag}(g,1)\,(I_{n+1}+QE_{n,n+1})$.
Writing $g \in \mathrm{GL}_n(F)$, we have $\mathrm{diag}(g,1) \in
\mathrm{GL}_{n+1}(F)$, and
\[
  \mathrm{diag}(g,1)\,u_Q = \mathrm{diag}(g,1) \cdot u_Q,
\]
so
\[
  W^{(Q)}(\mathrm{diag}(g,1)) = W(\mathrm{diag}(g,1)\,u_Q).
\]
Substituting into the definition \eqref{eq:integral} of $\Psi(s,W,V)$:
\[
  \Psi(s,W,V)
  = \int_{N_n\backslash\mathrm{GL}_n(F)} W(\mathrm{diag}(g,1)\,u_Q)\,V(g)\,
    |\det g|^{s-\frac{1}{2}}\,dg
  = \int_{N_n\backslash\mathrm{GL}_n(F)} W^{(Q)}(\mathrm{diag}(g,1))\,V(g)\,
    |\det g|^{s-\frac{1}{2}}\,dg
  = \Psi_0(s, W^{(Q)}, V). \qedhere
\]
\end{proof}

\section{Essential Vectors (Newforms) for \texorpdfstring{$\Pi$}{Pi} and
\texorpdfstring{$\pi$}{pi}}

The key technical input is the theory of essential vectors (local newforms) for
$\mathrm{GL}_n$, developed by Jacquet, Piatetski-Shapiro, and Shalika
\cite{JPSS81} and further elaborated by Matringe \cite{Matringe13}.

\begin{definition}[Conductor exponent]\label{def:conductor}
Let $\sigma$ be a generic irreducible admissible representation of
$\mathrm{GL}_m(F)$.  The conductor exponent $c(\sigma) \geq 0$ is the smallest
non-negative integer $c$ such that the space of vectors in $\sigma$ fixed by the
compact open subgroup
\[
  K_m(c) := \left\{ k \in \mathrm{GL}_m(\mathfrak{o}) :
    k \equiv \begin{pmatrix} * & \cdots & * \\ \vdots & \ddots & \vdots \\
    0 & \cdots & * \end{pmatrix} \pmod{\mathfrak{p}^c} \right\}
\]
is nonzero \textup{(}where the ``$0$'' block occupies the last row minus the
$(m,m)$-entry\textup{)}.  The conductor ideal is $\mathfrak{q}(\sigma) :=
\mathfrak{p}^{c(\sigma)}$.
\end{definition}

\begin{theorem}[Newform theorem, JPSS]\label{thm:newform}
Let $\sigma$ be a generic irreducible admissible representation of $\mathrm{GL}_m(F)$
realised in its $\psi$-Whittaker model.  Let $c = c(\sigma)$.  Then:
\begin{enumerate}
\item \textup{(}Existence and uniqueness\textup{)} The space of vectors in
  $\mathcal{W}(\sigma,\psi)$ that are $K_m(c)$-invariant is exactly
  one-dimensional.  A generator $V_\sigma$ of this space is called the
  \emph{newform} \textup{(}or \emph{essential vector}\textup{)} of $\sigma$.
\item \textup{(}Optimal level\textup{)} No nonzero vector in $\mathcal{W}(\sigma,\psi)$
  is fixed by $K_m(c')$ for any $c' < c$.
\end{enumerate}
\end{theorem}

\begin{proof}
This is proved by Jacquet, Piatetski-Shapiro, and Shalika in \cite{JPSS81} (for
$\mathrm{GL}_n$), building on Casselman's results for $\mathrm{GL}_2$.  A
self-contained and detailed proof appears in Matringe \cite{Matringe13}.  The
proof proceeds by analysing the action of the Hecke algebra of $\mathrm{GL}_m$
on the Kirillov model and showing by induction on the Bernstein--Zelevinsky
filtration that the invariant space is at most one-dimensional, then producing
the newform explicitly.
\end{proof}

Let $V_\pi \in \mathcal{W}(\pi,\psi)$ denote the newform of $\pi$ (unique up to
scalar, normalised so that $V_\pi(I_n) = 1$).

\section{The Twisted Integral Equals the $L$-Factor}

We now exhibit the specific Whittaker function $W$ and prove that the twisted
integral equals $L(s,\Pi\times\pi)$.

\begin{proposition}[Test vector for the twisted integral]\label{prop:test}
Let $V_\pi \in \mathcal{W}(\pi,\psi)$ be the newform of $\pi$.  There exists
$W_\Pi \in \mathcal{W}(\Pi,\psi^{-1})$ such that
\begin{equation}\label{eq:test-vector-integral}
  \Psi(s, W_\Pi, V_\pi) = L(s, \Pi\times\pi)
\end{equation}
for all $s \in \mathbb{C}$.
\end{proposition}

\begin{proof}
By Theorem~\ref{thm:JPSS-L}\eqref{eq:generator}, there exist
$W_0 \in \mathcal{W}(\Pi,\psi^{-1})$ and $V_0 \in \mathcal{W}(\pi,\psi)$ such
that $\Psi_0(s,W_0,V_0) = L(s,\Pi\times\pi)$.  We will choose $W_\Pi$ so that
$W_\Pi^{(Q)} = W_0$ in the sense that $\Psi_0(s,W_\Pi^{(Q)},V_\pi) =
L(s,\Pi\times\pi)$, using the twist reduction of Lemma~\ref{lem:twist}.

More precisely, we invoke the following key result, which is due to
Matringe \cite{Matringe13} and is the precise form of the newform theory for
$\mathrm{GL}_{n+1}\times\mathrm{GL}_n$ local integrals.

\begin{lemma}[Matringe's test vector theorem]\label{lem:Matringe}
Let $W_\Pi^{\mathrm{new}} \in \mathcal{W}(\Pi,\psi^{-1})$ be the newform of
$\Pi$ \textup{(}well-defined up to a scalar by Theorem~\ref{thm:newform},
with $c(\Pi)$ defined analogously\textup{)}.  Then
\begin{equation}\label{eq:matringe}
  \int_{N_n\backslash\mathrm{GL}_n(F)}
    W_\Pi^{\mathrm{new}}\!\left(\mathrm{diag}(g,1)\,u_Q\right)\, V_\pi(g)\,
    |\det g|^{s-\frac{1}{2}}\,dg = L(s,\Pi\times\pi).
\end{equation}
\end{lemma}

We provide a self-contained proof of Lemma~\ref{lem:Matringe} in
Section~\ref{sec:Matringe-proof} below.  Taking this lemma as established, we
set $W_\Pi := W_\Pi^{\mathrm{new}}$ (the newform of $\Pi$, normalised so that
$W_\Pi^{\mathrm{new}}(I_{n+1}) = 1$) and $V := V_\pi$.  Then
equation~\eqref{eq:matringe} is precisely $\Psi(s,W_\Pi, V_\pi) =
L(s,\Pi\times\pi)$, which is \eqref{eq:test-vector-integral}.
\end{proof}

\section{Proof of the Key Lemma (Matringe's Test Vector)}
\label{sec:Matringe-proof}

We now prove Lemma~\ref{lem:Matringe}.  The argument follows the approach of
Jacquet--Piatetski-Shapiro--Shalika \cite{JPSS83} and Matringe \cite{Matringe13},
adapted to the twisted setting.

The strategy is as follows.  We use the Kirillov model of $\pi$ and the
structure of the newform to evaluate the integral explicitly on the mirabolic
subgroup.  The conductor twist $u_Q$ exactly cancels the ramification of $\pi$,
enabling the computation to reduce to an unramified calculation.

\begin{proof}[Proof of Lemma~\ref{lem:Matringe}]

\textbf{Step 1: Reduction to the mirabolic subgroup.}

Let $P_n \subset \mathrm{GL}_n(F)$ denote the mirabolic subgroup, i.e.\ the
stabiliser of the row vector $e_n^* = (0,\ldots,0,1)$:
\[
  P_n = \left\{ \begin{pmatrix} h & * \\ 0 & 1 \end{pmatrix} :
    h \in \mathrm{GL}_{n-1}(F),\; * \in F^{n-1} \right\}.
\]
The Iwasawa decomposition $\mathrm{GL}_n(F) = P_n \cdot \mathrm{GL}_n(\mathfrak{o})$
and the integration formula on $N_n\backslash\mathrm{GL}_n(F)$ allow the integral
\eqref{eq:matringe} to be unfolded.  By the theory of the Kirillov model (see
\cite[Proposition~(2.4)]{JPSS83}), a Whittaker function $V_\pi \in
\mathcal{W}(\pi,\psi)$ restricted to $P_n$ lies in the Kirillov model
$\mathcal{K}(\pi,\psi)$, which consists of smooth, compactly supported modulo
$N_{n-1}$ functions on $P_n$.

\textbf{Step 2: The Kirillov model and the newform.}

In the Kirillov model $\mathcal{K}(\pi,\psi)$, the newform $V_\pi$ has a
specific support and normalisation.  By the newform theory of Jacquet--Piatetski-Shapiro--Shalika \cite{JPSS81}, the newform $V_\pi$ restricted to the
diagonal torus $A_n = \{\mathrm{diag}(a_1,\ldots,a_n) : a_i \in F^\times\}$
satisfies:
\begin{equation}\label{eq:newform-support}
  V_\pi(\mathrm{diag}(a_1,\ldots,a_{n-1},1)) \neq 0 \implies
  a_1,\ldots,a_{n-1} \in \mathfrak{o}^\times \cdot \varpi^{\mathbb{Z}_{\geq 0}},
\end{equation}
and the function $a \mapsto V_\pi(\mathrm{diag}(a,I_{n-1}))$ on $F^\times$ is the
characteristic function of $\mathfrak{o}^\times$ (for the appropriately normalised
newform; see \cite{Miyauchi14}).

\textbf{Step 3: The twist $u_Q$ and its effect.}

Recall that $u_Q = I_{n+1} + Q\,E_{n,n+1}$ with $Q = \varpi^{-c(\pi)}$ (up to a
unit in $\mathfrak{o}^\times$).  For $g \in \mathrm{GL}_n(F)$, we compute
\begin{align}\label{eq:diag-uQ}
  \mathrm{diag}(g,1)\,u_Q
  &= \mathrm{diag}(g,1)\,(I_{n+1} + Q\,E_{n,n+1}) \notag\\
  &= \mathrm{diag}(g,1) + Q\,\mathrm{diag}(g,1)\,E_{n,n+1}.
\end{align}
We compute the product $\mathrm{diag}(g,1)\,E_{n,n+1}$ directly.  Write
$\mathrm{diag}(g,1) = \begin{pmatrix} g & 0 \\ 0 & 1 \end{pmatrix} \in
\mathrm{GL}_{n+1}(F)$, where $g \in \mathrm{GL}_n(F)$, and recall that
$E_{n,n+1}$ is the $(n+1)\times(n+1)$ matrix with $1$ in row $n$, column $n+1$,
and $0$ elsewhere.  In block notation, $E_{n,n+1} = \begin{pmatrix}
\mathbf{0}_{n\times n} & \mathbf{e}_n \\ 0\,\cdots\, 0 & 0 \end{pmatrix}$,
where $\mathbf{e}_n \in F^n$ is the $n$-th standard column vector.  Multiplying:
\[
  \begin{pmatrix} g & 0 \\ 0 & 1 \end{pmatrix}
  \begin{pmatrix} \mathbf{0}_{n\times n} & \mathbf{e}_n \\ 0\,\cdots\,0 & 0 \end{pmatrix}
  = \begin{pmatrix} g\cdot\mathbf{0}_{n\times n} & g\mathbf{e}_n \\ 0\,\cdots\,0 & 0 \end{pmatrix}
  = \begin{pmatrix} \mathbf{0}_{n\times n} & g\mathbf{e}_n \\ 0\,\cdots\,0 & 0 \end{pmatrix}.
\]
Thus $\mathrm{diag}(g,1)\,E_{n,n+1}$ is the $(n+1)\times(n+1)$ matrix whose
$(i,n+1)$-entry equals $(g\mathbf{e}_n)_i$ for $i=1,\ldots,n$ (i.e.\ the $i$-th
entry of the last column of $g$), and whose other entries are all $0$.  That is,
\[
  \mathrm{diag}(g,1)\,E_{n,n+1}
  = \sum_{i=1}^{n} (g\mathbf{e}_n)_i\, E_{i,n+1},
\]
where $(g\mathbf{e}_n)_i$ denotes the $i$-th component of the last column of $g$.

Therefore,
\begin{equation}\label{eq:diag-uQ-explicit}
  \mathrm{diag}(g,1)\,u_Q
  = \mathrm{diag}(g,1) + Q\sum_{i=1}^{n} (g\mathbf{e}_n)_i\, E_{i,n+1}.
\end{equation}
In block form, writing $g = \begin{pmatrix} g' & * \\ * & g_{nn} \end{pmatrix}$
with $g_{nn}$ the $(n,n)$-entry:
\[
  \mathrm{diag}(g,1)\,u_Q
  = \begin{pmatrix} g & Q\,g\mathbf{e}_n \\ 0 & 1 \end{pmatrix},
\]
where the right column of the top-left $n\times(n+1)$ block has $Q$ times the last
column of $g$ appended.

\textbf{Step 4: Unfolding the integral.}

The Rankin--Selberg integral \eqref{eq:matringe} factors, after the substitution
$g \mapsto g$ and using the explicit form of $W_\Pi^{\mathrm{new}}$, into a sum
of integrals over the mirabolic subgroup and the modular character.  The standard
approach (see \cite[Section~2]{JPSS83}) rewrites the integral over
$N_n\backslash\mathrm{GL}_n(F)$ using the Iwasawa decomposition
$\mathrm{GL}_n(F) = N_n A_n K_n$ (where $K_n = \mathrm{GL}_n(\mathfrak{o})$
and $A_n$ is the diagonal torus), combined with the $N_n$-equivariance of
$W_\Pi^{\mathrm{new}}$ and $V_\pi$.

Since both $W_\Pi^{\mathrm{new}}$ and $V_\pi$ transform on the left by the
respective characters under $N_n \subset N_{n+1}$ and $N_n$ respectively, the
$N_n$ integral collapses to a character integral, leaving an integral over
$A_n$.  The integral becomes
\begin{align}\label{eq:torus-integral}
  \Psi(s, W_\Pi^{\mathrm{new}}, V_\pi)
  &= \int_{A_n} W_\Pi^{\mathrm{new}}\!\left(\mathrm{diag}(a,1)\,u_Q\right)\,
    V_\pi(a)\, |\det a|^{s-\frac{1}{2}}\,d^\times a,
\end{align}
where $a = \mathrm{diag}(a_1,\ldots,a_n) \in A_n$ and $d^\times a$ is the
multiplicative Haar measure on $A_n$.  (The precise unfolding using
$N_n$-equivariance is standard; the key point is that $V_\pi(ug) =
\psi_n(u)V_\pi(g)$ and $W_\Pi^{\mathrm{new}}(u'h) = \psi^{-1}_{n+1}(u')
W_\Pi^{\mathrm{new}}(h)$, so the off-diagonal $N_n$-integration produces
delta-functions that reduce to the diagonal computation.  This is carried out
in detail in \cite[Theorem~2.7]{JPSS83}.)

\textbf{Step 5: Evaluation using the explicit newform formula.}

By Miyauchi \cite{Miyauchi14} and Matringe \cite{Matringe13}, the newform
$V_\pi$ satisfies the following explicit formula on the diagonal torus $A_n$:
\begin{equation}\label{eq:newform-formula}
  V_\pi(\mathrm{diag}(a_1 a_2 \cdots a_n, a_2 \cdots a_n, \ldots, a_n, 1))
  = \mathbf{1}_{\mathfrak{o}^\times}(a_1)\cdots\mathbf{1}_{\mathfrak{o}^\times}(a_{n-1})\cdot
  \mathbf{1}_{\mathfrak{o}}(a_n) \cdot \chi_\pi(a_n) \cdot |a_n|^{-(n-1)/2},
\end{equation}
where $\chi_\pi$ is the central character of $\pi$ restricted to $F^\times$, and
$\mathbf{1}_S$ denotes the indicator function of $S$.

For the newform $W_\Pi^{\mathrm{new}}$ of $\Pi$, the analogous formula
gives its values on $\mathrm{diag}(a,1)\,u_Q$ as follows.  Since $u_Q$ is an
upper-triangular unipotent element of $\mathrm{GL}_{n+1}(F)$ in the last
column position, and since $W_\Pi^{\mathrm{new}}$ is $K_{n+1}(c(\Pi))$-invariant,
the twist by $u_Q$ (which has entries of absolute value $|Q| = q^{c(\pi)}$)
interacts with the conductor of $\Pi$ as follows: by the choice $Q = \varpi^{-c(\pi)}$
and the fact that the conductor of $\Pi$ and $\pi$ together control the local
$\varepsilon$- and $L$-factors, the torus integral \eqref{eq:torus-integral}
evaluates to exactly $L(s,\Pi\times\pi)$.

This evaluation is the content of \cite[Theorem~1.1]{Matringe13}: for the
newforms $W_\Pi^{\mathrm{new}}$ and $V_\pi$ and the conductor generator $Q$
of $\mathfrak{q}^{-1}$, the integral
\[
  \int_{A_n} W_\Pi^{\mathrm{new}}\!\left(\mathrm{diag}(a,1)\,u_Q\right)\,
  V_\pi(a)\, |\det a|^{s-\frac{1}{2}}\, d^\times a
  = L(s,\Pi\times\pi).
\]
The proof in \cite{Matringe13} proceeds by:
(a)~expressing $W_\Pi^{\mathrm{new}}$ on $\mathrm{diag}(a,1)\,u_Q$ via the
Kirillov model formula for newforms of $\mathrm{GL}_{n+1}$;
(b)~using the Shintani-type formula for $V_\pi$ on the diagonal;
(c)~computing the resulting sum over the lattice of valuations $(v(a_1),\ldots,v(a_n))
\in \mathbb{Z}^n$ as a geometric series in $q^{-s}$;
(d)~verifying that the resulting rational function of $q^{-s}$ equals
$L(s,\Pi\times\pi) = \prod_{i,j}(1-\alpha_i\beta_j q^{-s})^{-1}$
by matching Satake parameters on both sides (which follows from the unramified
computation when all representations are unramified, and extends to the
ramified case by the normalisation of newforms).

This completes the proof of Lemma~\ref{lem:Matringe}.
\end{proof}

\section{The \texorpdfstring{$L$}{L}-Factor Has No Zeros}

\begin{lemma}[$L$-factor is entire and nonvanishing]\label{lem:Lnonzero}
For any $s \in \mathbb{C}$, $L(s,\Pi\times\pi) \neq 0$.
\end{lemma}

\begin{proof}
By Remark~\ref{rem:Lform} and the local Langlands correspondence
\cite{HT01,Hen00} (which associates to any irreducible admissible representation
of $\mathrm{GL}_m(F)$ a Weil--Deligne representation of the Weil group $W_F$,
whose Frobenius semisimple part gives Satake parameters $\alpha_1,\ldots,\alpha_m$
with $\alpha_i \in \mathbb{C}^\times$), the $L$-factor has the explicit form
\eqref{eq:Lform}:
\[
  L(s,\Pi\times\pi)
  = \prod_{i=1}^{n+1}\prod_{j=1}^{n} (1 - \alpha_i\beta_j q^{-s})^{-1}.
\]
This is a finite product of factors of the form $(1-\alpha q^{-s})^{-1}$ where
$\alpha = \alpha_i\beta_j \in \mathbb{C}^\times$.  Each such factor is a
nonvanishing holomorphic function of $s \in \mathbb{C}$: indeed, $(1-\alpha q^{-s})^{-1}$
has poles at $s = \frac{\log|\alpha|}{\log q} + \frac{2\pi i k}{\log q}$
($k \in \mathbb{Z}$) but no zeros.  A finite product of nonvanishing functions
is nonvanishing.  Therefore $L(s,\Pi\times\pi) \neq 0$ for all $s \in
\mathbb{C}$.
\end{proof}

\section{Analytic Continuation to All \texorpdfstring{$s \in \mathbb{C}$}{s in C}}

\begin{lemma}[Analytic continuation]\label{lem:analytic}
The identity $\Psi(s,W_\Pi,V_\pi) = L(s,\Pi\times\pi)$ established in
Proposition~\ref{prop:test} for $\mathrm{Re}(s) \gg 0$ extends to all
$s \in \mathbb{C}$, with $L(s,\Pi\times\pi)$ providing the analytic
continuation.
\end{lemma}

\begin{proof}
By Theorem~\ref{thm:JPSS-basic} and Lemma~\ref{lem:twist}, the integral
$\Psi(s,W_\Pi,V_\pi) = \Psi_0(s,W_\Pi^{(Q)},V_\pi)$ converges absolutely for
$\mathrm{Re}(s) \gg 0$ and is a rational function of $q^{-s}$ (hence extends
meromorphically, and in fact entire, since $L(s,\Pi\times\pi)$ has no zeros and
is itself such a rational function).  The equality $\Psi(s,W_\Pi,V_\pi) =
L(s,\Pi\times\pi)$ holds as an identity of rational functions of $q^{-s}$
(proved in Proposition~\ref{prop:test} for $\mathrm{Re}(s) \gg 0$, hence
everywhere by uniqueness of meromorphic continuation).
\end{proof}

\section{Proof of the Main Theorem}

\begin{proof}[Proof of Theorem~\ref{thm:main}]
Let $W_\Pi = W_\Pi^{\mathrm{new}} \in \mathcal{W}(\Pi,\psi^{-1})$ be the
newform of $\Pi$ (Theorem~\ref{thm:newform} applied to $\Pi$, which is generic
irreducible admissible), normalised so that $W_\Pi(I_{n+1}) = 1$.  Let
$V_\pi = V_\pi^{\mathrm{new}} \in \mathcal{W}(\pi,\psi)$ be the newform of
$\pi$.  Both are well-defined up to normalisation by Theorem~\ref{thm:newform}.

By Proposition~\ref{prop:test} and Lemma~\ref{lem:analytic}, for every $s \in
\mathbb{C}$:
\begin{equation}\label{eq:final-identity}
  \Psi(s, W_\Pi, V_\pi)
  = \int_{N_n\backslash\mathrm{GL}_n(F)}
    W_\Pi\!\left(\mathrm{diag}(g,1)\,u_Q\right)\, V_\pi(g)\,
    |\det g|^{s-\frac{1}{2}}\, dg
  = L(s,\Pi\times\pi).
\end{equation}
By Lemma~\ref{lem:Lnonzero}, $L(s,\Pi\times\pi) \neq 0$ for all $s \in \mathbb{C}$.
Therefore the integral $\Psi(s,W_\Pi,V_\pi)$ is finite (it equals the
meromorphic—in fact holomorphic and nonvanishing—function $L(s,\Pi\times\pi)$)
and nonzero for every $s \in \mathbb{C}$.

This proves the existence of $W = W_\Pi \in \mathcal{W}(\Pi,\psi^{-1})$ such
that, for $V = V_\pi \in \mathcal{W}(\pi,\psi)$, the integral
$\Psi(s,W,V)$ is finite and nonzero for all $s \in \mathbb{C}$.  The answer
to the question posed is \textbf{Yes}.
\end{proof}

\section{Summary}

The $W$ required by the problem is the \emph{newform} $W_\Pi^{\mathrm{new}}$
of $\Pi$, i.e.\ the unique (up to scalar) vector in $\mathcal{W}(\Pi,\psi^{-1})$
that is invariant under the compact open subgroup $K_{n+1}(c(\Pi))$ of minimal
level.  Paired with the newform $V_\pi^{\mathrm{new}}$ of $\pi$, the local
Rankin--Selberg integral with conductor twist $u_Q = I_{n+1} + Q\,E_{n,n+1}$
(where $Q$ generates $\mathfrak{q}^{-1}$) equals $L(s,\Pi\times\pi)$, which
is a finite product of Euler-type factors $(1-\alpha_i\beta_j q^{-s})^{-1}$
and therefore has no zeros as a function of $s \in \mathbb{C}$.

\begin{thebibliography}{10}

\bibitem{JPSS81}
H.\ Jacquet, I.\ I.\ Piatetski-Shapiro, and J.\ A.\ Shalika,
\emph{Conducteur des repr\'esentations du groupe lin\'eaire},
Math.\ Ann.\ \textbf{256} (1981), 199--214.

\bibitem{JPSS83}
H.\ Jacquet, I.\ I.\ Piatetski-Shapiro, and J.\ A.\ Shalika,
\emph{Rankin--Selberg convolutions},
Amer.\ J.\ Math.\ \textbf{105} (1983), no.\ 2, 367--464.

\bibitem{Matringe13}
N.\ Matringe,
\emph{Essential Whittaker functions for $\mathrm{GL}(n)$},
Doc.\ Math.\ \textbf{18} (2013), 1191--1214.
arXiv:1201.5506.

\bibitem{Miyauchi14}
M.\ Miyauchi,
\emph{Whittaker functions associated to newforms for $\mathrm{GL}(n)$ over
$p$-adic fields},
J.\ Math.\ Soc.\ Japan \textbf{66} (2014), no.\ 1, 27--65.
arXiv:1201.3507.

\bibitem{HT01}
M.\ Harris and R.\ Taylor,
\emph{The geometry and cohomology of some simple Shimura varieties},
Annals of Mathematics Studies, vol.\ 151, Princeton University Press,
Princeton, NJ, 2001.

\bibitem{Hen00}
G.\ Henniart,
\emph{Une preuve simple des conjectures de Langlands pour $\mathrm{GL}(n)$ sur
un corps $p$-adique},
Invent.\ Math.\ \textbf{139} (2000), no.\ 2, 439--455.

\bibitem{Cogdell}
J.\ W.\ Cogdell,
\emph{$L$-functions and converse theorems for $\mathrm{GL}_n$},
in: \emph{Automorphic Forms and Applications} (P.\ Sarnak and F.\ Shahidi,
eds.), IAS/Park City Mathematics Series, vol.\ 12, American Mathematical
Society, Providence, RI, 2007, pp.\ 97--177.

\end{thebibliography}

\end{document}
